\chapter{Metodología}
\section{Diseño y alcance}
Se realizó un estudio computacional de un monopolo inmerso en cuatro medios homogéneos a 3 GHz. Se resolvieron las ecuaciones de Maxwell mediante FDTD en Meep, con coordenadas cilíndricas y modo azimutal $m=0$. La evidencia se organizó en diez configuraciones históricas de atenuación y cinco corridas adicionales de balance local. Los archivos y sus versiones se identificaron en el Anexo~\ref{an:fdtd-fuentes}.

La penetración se examinó mediante el decaimiento dentro del material. La transmisión a través de interfaces entre materiales y la pérdida de inserción de una pared quedaron fuera del dominio simulado. Esta delimitación se utilizó al evaluar los objetivos de la tesis.

Se trabajó en un equipo con procesador AMD Ryzen 5 4600H. La campaña de balance se ejecutó desde VS Code conectado a Ubuntu bajo WSL2, con Python 3.13.13, Meep 1.33.0-beta y un hilo de cómputo. Los tiempos y la memoria máxima se registraron por corrida.

\section{Modelo físico y geometría}
Se fijaron $f=3.0\times10^9$ Hz, $c=299792458$ m/s y $\lambda_0=c/f$. Las longitudes se introdujeron en milímetros.
Los medios analizados fueron definidos mediante su permitividad relativa $\varepsilon_r$ y su conductividad $\sigma$.

\begin{table}[H]
\centering
\begin{tabular}{|l|c|c|}
\hline
\textbf{Medio} & $\boldsymbol{\varepsilon_r}$ & $\boldsymbol{\sigma}$ (\textbf{S/m}) \\ \hline
Vacío & 1.0 & 0.0 \\ \hline
Concreto & 4.5 & 0.05 \\ \hline
Agua salada & 78.0 & 4.0 \\ \hline
Tejido muscular & 52.7 & 2.2 \\ \hline
\end{tabular}
\caption{Propiedades electromagnéticas utilizadas para definir los medios simulados a la frecuencia de trabajo.}
\label{tab:materiales_metodologia}
\end{table}

Cada material fue simulado por separado. El medio activo se seleccionó mediante la variable \texttt{nombre\_medio}. A partir de esta selección, el programa recuperó automáticamente los valores de $\varepsilon_r$ y $\sigma$.

Con estos parámetros se calculó la tangente de pérdidas,

\begin{equation}
\tan\delta=
\frac{\sigma}
{\omega\varepsilon_0\varepsilon_r},
\end{equation}

además de las constantes de atenuación $\alpha$ y fase $\beta$.

La longitud de onda dentro del medio fue calculada mediante

\begin{equation}
\lambda_m=\frac{2\pi}{\beta}.
\end{equation}

Estos valores fueron obtenidos antes de iniciar la simulación FDTD y se utilizaron para adaptar las dimensiones del dominio a cada material.

La conductividad física expresada en $\mathrm{S/m}$ fue convertida a las unidades utilizadas internamente por MEEP. Después de esta conversión, cada medio se definió mediante \texttt{mp.Medium}, usando su permitividad relativa y su conductividad normalizada.

La geometría estuvo formada por un plano conductor perfecto y un monopolo cilíndrico. El radio del monopolo fue de $0.2$ mm, por lo que su diámetro fue de $0.4$ mm.

El programa definió inicialmente

\begin{equation}
L_{\mathrm{antena}}=\frac{\lambda_m}{2}.
\end{equation}

La altura del brazo metálico se calculó como

\begin{equation}
h_{\mathrm{brazo}}
=
\frac{L_{\mathrm{antena}}-g}{2},
\end{equation}

donde la separación de alimentación inicial fue

\begin{equation}
g=0.4\ \mathrm{mm}.
\end{equation}

Debido a esta definición, la longitud física del monopolo cambió con el material, ya que $\lambda_m$ también cambió. Por lo tanto, las simulaciones compararon monopolos ajustados a la longitud de onda de cada medio y no una única antena con dimensiones físicas constantes.

El dominio radial útil fue definido como

\begin{equation}
R_{\max}=5\lambda_m.
\end{equation}

La región absorbente tuvo un espesor de

\begin{equation}
d_{\mathrm{abs}}=2\lambda_m.
\end{equation}

Para valores de $\tan\delta\leq1$, el programa utilizó capas PML en las direcciones radial y axial. También se dejó disponible una condición alternativa con \texttt{Absorber} para valores de $\tan\delta>1$.


\section{Resoluciones y extracción de atenuación}
Se compararon los cuatro medios a 5 px/mm, equivalentes a $h=0.2$ mm y dos celdas sobre el diámetro del conductor. Las resoluciones históricas fueron 5 px/mm en vacío, 5 y 10 en concreto, 5, 10, 20 y 40 en agua salada, y 5, 10 y 20 en músculo. Cada resolución se identificó a partir del archivo de la corrida. En el programa adicional de balance se seleccionó mediante el argumento \texttt{--resolucion}.

El simulador histórico extrajo $E_z$ y $H_\phi$ sobre $z=0$ al concluir el avance temporal. Se calculó $S_r=-\operatorname{Re}(E_zH_\phi^*)/2$ y se ajustó el logaritmo del perfil $|S_r|r^2$. El módulo eliminó el signo del flujo de esta serie radial. Los residuos se definieron como observación menos ajuste.

Se emplearon ventanas entre $1.5\lambda_m$ y $4\lambda_m$ en vacío y concreto, y entre $2\lambda_m$ y $4\lambda_m$ en agua salada y músculo, con puntos superiores a $-80$ dB. El modelo M0 ajustó una recta, M1 añadió una sinusoide y M2 añadió un término logarítmico. Se restringió la atenuación a valores no negativos y se prefirió M0, salvo una mejora del AIC superior a diez unidades. Se exigieron $R^2>0.85$ y RMSE inferior a 0.5 dB para considerar utilizable la extracción.

Los IC del 95\,\% se obtuvieron de los percentiles 2.5 y 97.5 de 1500 intentos de remuestreo por bloques de residuos, con semilla cero. Se reprodujo el ajuste en los cinco perfiles numéricos disponibles y se documentó la corrección del índice de inicio de los bloques. Para las otras configuraciones se recuperaron los IC de los informes de las figuras. El procedimiento y sus limitaciones se detallaron en el Anexo~\ref{an:fdtd-fuentes}.

\section{Medición del balance local}
Se conservaron la fuente continua, los parámetros del medio y la geometría del simulador histórico. Se definió un anillo homogéneo sin fuente ni conductor, con $r\in[\lambda_m,2\lambda_m]$ y $z\in[0.25\lambda_m,1.25\lambda_m]$. Se midieron las cuatro caras y la pérdida volumétrica de Joule mediante las Ecuaciones~\ref{eq:joule-anular} y~\ref{eq:balance-normalizado}.

En unidades temporales de Meep, se utilizó $t_{\rm vuelo}=R_{\max}\sqrt{\varepsilon_r}$ y
\begin{equation}
 t_{\rm est}=\begin{cases}3t_{\rm vuelo}+20\lambda_0,&\tan\delta<0.1,\\1.5t_{\rm vuelo}+10\lambda_0,&\tan\delta\geq0.1.\end{cases}
\end{equation}
Se tomó una muestra al terminar este intervalo y otra aproximadamente un período después. Antes de cada medición se sincronizaron los campos eléctrico y magnético. Se interpolaron sus valores en puntos físicos del anillo y se aplicó la regla trapezoidal con el factor cilíndrico $2\pi r$.

Se calcularon los flujos netos por cara y se agruparon por su signo para obtener entrada y salida. Se verificó el residuo $R=P_{\rm ent}-P_{\rm sal}-P_{\rm J}$ y el cambio temporal de cada término. Se dispuso de balances a 5 px/mm para los cuatro medios y de un refinamiento a 10 px/mm en agua salada. Las pruebas breves de instrumentación se excluyeron de los resultados físicos.

\section{Unidades y normalización}\label{sec:unidades-campana}
Se adoptó la longitud de referencia $a=1$ mm. Las coordenadas numéricas se relacionaron con las físicas mediante $x_{\rm SI}=a x_{\rm M}$ y el tiempo mediante $t_{\rm SI}=(a/c)t_{\rm M}$. Así, una unidad temporal correspondió a 3.33564 ps y $f_{\rm M}=fa/c=0.0100069$. El período de la fuente fue de aproximadamente 99.9308 unidades temporales, equivalente a 333.333 ps.

La conductividad introducida en Meep se calculó como
\begin{equation}
 \sigma_D=\frac{\sigma_{\rm SI}a}{\varepsilon_0\varepsilon_r c}.
\end{equation}
Este valor normalizado multiplicó a $D$ en la ecuación de evolución. Por ello, la densidad de pérdida usada en la integración fue $q_{\rm M}=\varepsilon_r\sigma_D|E_{\rm M}|^2/2$.

Para expresar los campos y las potencias en SI se requirió además una escala de campo $E_*$, en V/m. Con $Z_0=\sqrt{\mu_0/\varepsilon_0}$, se obtuvo
\begin{equation}
 E_{\rm SI}=E_*E_{\rm M},\quad H_{\rm SI}=\frac{E_*}{Z_0}H_{\rm M},\quad P_{\rm SI}=\frac{E_*^2a^2}{Z_0}P_{\rm M}.
\end{equation}
La excitación utilizada dejó $E_*$ sin calibrar. Los valores de potencia se informaron como $P/P_*$, con $P_*=E_*^2a^2/Z_0$, y las fracciones de potencia fueron adimensionales. El tiempo de pared en segundos describió el costo de ejecución, mientras que $t_{\rm M}$ describió la evolución electromagnética.

\section{Protocolo de comparación y criterios de interpretación}
Se propuso un procedimiento reproducible que reunió la identificación del material y la malla, la extracción radial con ventana declarada, el informe de IC y residuos, la comparación teórica y la comprobación independiente del balance. Para comparar resoluciones del balance se conservaron los límites físicos del mismo anillo.

La sensibilidad de una magnitud $Q$ se expresó como $\delta_Q=100|Q_5-Q_{\rm fina}|/|Q_{\rm fina}|$. Se examinó el umbral de 5\,\% propuesto en el diseño como criterio interno de comparación. Su cumplimiento se evaluó por magnitud. La fracción disipada y las potencias absolutas recibieron evaluaciones independientes. Cuando la referencia fue cero se utilizó una diferencia absoluta.

La discrepancia teórica de atenuación, el IC condicional del ajuste, el cambio temporal y la diferencia entre mallas se reportaron individualmente. Se evaluó el orden aparente en agua salada, atendiendo al cambio de modelo de ajuste. La extrapolación al continuo quedó condicionada a la identificación de un régimen asintótico estable.

Se conservaron los perfiles originales, las figuras y los informes históricos. Cada nueva corrida de balance produjo un JSON fechado, con dos muestras, dominio, versión, tiempo de ejecución y memoria máxima. Los CSV derivados se vincularon con ese registro mediante un índice y hashes. La organización se documentó en el Anexo~\ref{an:balance-registros}.
